% muLLM: Cost-Optimal LLM Routing via Lightweight Intent Classification
% Compact anonymized version — 4-page workshop submission
% Compile with: pdflatex mullm_paper_compact_anonymized && bibtex mullm_paper_compact_anonymized
%               pdflatex mullm_paper_compact_anonymized && pdflatex mullm_paper_compact_anonymized

\documentclass[11pt]{article}

\usepackage[review]{acl}

\usepackage{times}
\usepackage{latexsym}
\usepackage[T1]{fontenc}
\usepackage[utf8]{inputenc}
\usepackage{microtype}
\usepackage{inconsolata}
\usepackage{graphicx}
\usepackage{booktabs}
\usepackage{listings}
\usepackage{amsmath}

\lstset{
  basicstyle=\small\ttfamily,
  breaklines=true,
  frame=none,
  columns=fullflexible,
  keepspaces=true,
}

\title{muLLM: Cost-Optimal LLM Routing via Lightweight Intent Classification}

\author{Anonymous Submission}

\begin{document}
\maketitle

%% ============================================================
\begin{abstract}
Most LLM queries don't need a frontier model. muLLM is a four-tier
routing system built on one empirical claim: \textbf{the majority of
real production queries can be answered at zero marginal cost without
quality loss}. The system routes each query to the cheapest tier that
can handle it: a deterministic lookup table ($<$1\,ms), semantic vector
cache ($\sim$14\,ms), local 30B model ($\sim$5\,s), and cloud fallback.
Over 15,363 queries across 22 days, muLLM spent \$27.87 against a \$772
blended cloud baseline --- a \textbf{96.4\% cost reduction}; 87.5\% of
queries were resolved free. The local tier achieves 100\% pass@1 on
HumanEval at \$0.00. A 44\,MB DeBERTa-v3-small classifier achieves AIQ
0.6673 on RouterBench, exceeding all published values. The key emergent
property: costs fall as usage grows.
\end{abstract}

%% ============================================================
\section{Introduction}

Deploying LLMs in production carries a fundamental tension: capable
frontier models are expensive at scale, while cheaper alternatives
sacrifice quality. Existing routers
(FrugalGPT~\cite{chen2023frugalgpt}, RouteLLM~\cite{ong2024routellm})
use binary or cascade routing --- choosing between a cheap and an
expensive model. This leaves two cost-free tiers on the table:
deterministic pattern matching and semantic caching, which together
resolve the majority of real production queries before any model is
invoked.

\textbf{muLLM} introduces a four-tier hierarchy that changes the cost
structure of LLM deployment. Over a 22-day production window spanning
15,363 organic queries (software engineering, research, writing, data
analysis), the system achieved 96.4\% cost reduction at production
quality. The 12.5\% cloud escalation rate is inflated by three
dedicated benchmark days; on organic workdays, the free-tier rate
exceeds 91\%.

Four contributions distinguish this work: (1) \textbf{Tier~0 groundtruth
short-circuiting} --- 1,313 deterministic patterns resolving 11.8\% of
queries in $<$1\,ms with zero hallucination risk; (2) \textbf{Bash-as-Truth
and Bash-as-Judge} --- evaluation primitives that eliminate LLM
hallucination at the resolver and judge layers; (3) \textbf{ELO-cache
semantics} --- a cache that converges toward the best-ever answer for
each semantic neighborhood; (4) \textbf{m$\mu$PETS} --- a routing-specific
Pareto efficiency metric capturing cost, accuracy, and latency jointly
(see Appendix).

%% ============================================================
\section{System Architecture}

\begin{lstlisting}
User Query
    |
    v
[Tier 0: Groundtruth LUT]  --- <1ms . $0.00 . 11.8%
    | no match
    v
[Tier 1: Semantic Cache]   --- ~14ms . $0.00 . 27.1%
    | cosine < 0.98
    v
[DeBERTa Classifier]       --- 44MB CPU-only . 82.4%
    |
    +-- complexity 1-2 -> [Tier 2a: Local 30B ~5s . $0.00]
    +-- complexity 3-5 -> [Tier 2b: Local 30B ~8s . $0.00]
                                   |
                         [Quality Gate] --- pass -> respond
                                   | fail
                                   v
                         [Tier 3: Cloud API . 2-3s . $0.001-$0.15]
                                   |
                         [Update Cache + Score Log]
\end{lstlisting}

\textbf{Tier~0 --- Groundtruth LUT} ($<$1\,ms, \$0.00). 1,313+
deterministic patterns across 13~categories: arithmetic, capitals, HTTP
status codes, git commands, Big-O complexity, date/time, unit
conversion, periodic table, ASCII, regex. A Python \texttt{Counter}
answers ``How many r's in strawberry?'' in 4\,ms with zero hallucination
risk. Categories derived by analyzing the first 10,284 production
queries for recurring deterministic-answer patterns.

\textbf{Tier~1 --- Semantic Cache} ($\sim$14\,ms, \$0.00). ChromaDB
vector store with cosine similarity threshold 0.98. ELO-cache semantics:
entries are only displaced by demonstrably better answers from
escalation events, causing the cache to converge toward best-ever
answers over time. Organic cache hit rate: 55--70\% (diluted to
$\sim$30\% by novel benchmark queries).

\textbf{Tier~2 --- Local Inference} ($\sim$5\,s, \$0.00). Ollama
backend with qwen3-coder:30b (MoE, $\sim$3B active parameters, $\sim$20\,GB
VRAM at Q4\_K\_M). DeBERTa-v3-small classifies intent (CODE, RESEARCH,
CREATIVE, CONVERSATION, DEPLOY, NOTE, LOOKUP, VISION) and complexity
(1--5). CPU-only inference, 14\,ms latency, 44\,MB model size.

\textbf{Tier~3 --- Cloud API} ($\sim$2--3\,s, \$0.001--\$0.15/query).
Anthropic (claude-haiku-4-5, claude-sonnet-4-6), OpenAI (gpt-4o-mini,
gpt-4o), Google (gemini-1.5-flash, gemini-1.5-pro). Hard spend limits
enforced at \$7/\$10/\$50 session thresholds. OpenAI-compatible
\texttt{/v1/chat/completions} for drop-in compatibility.

The routing decision (Tier~0 + classifier) adds 23\,$\mu$s median
overhead --- below perceptible latency.

%% ============================================================
\section{Evaluation}

\subsection{Cost Reduction (Production, 22-Day Window)}

\begin{table}[t]
\centering
\small
\begin{tabular}{lc}
\toprule
\textbf{Metric} & \textbf{Value} \\
\midrule
Total queries & 15,363 \\
Actual spend & \$27.87 \\
Blended cloud baseline & \$772 \\
Cost reduction & \textbf{96.4\%} \\
Token cost avoided & \textbf{99.7\%} \\
Free-tier query rate & \textbf{87.5\%} \\
\bottomrule
\end{tabular}
\caption{Production cost figures. Baseline: Sonnet pricing for routine queries, Opus pricing for complex queries.}
\label{tab:cost}
\end{table}

\subsection{Benchmark Results}

\begin{table}[t]
\centering
\small
\begin{tabular}{lccc}
\toprule
\textbf{Benchmark} & \textbf{Score} & \textbf{Cost} & \textbf{vs.\ Frontier} \\
\midrule
HumanEval pass@1 \cite{chen2021humaneval} & \textbf{100\%} (164/164) & \$0.00 & +8pp vs GPT-4o \\
MultiPL-E 6-lang \cite{cassano2022multipl} & \textbf{98.3\%} (118/120) & \$0.00 & parity GPT-4 \\
MultiPL-E 17-lang & \textbf{92.2\%} (249/270) & \$0.00 & strong \\
MMLU \cite{hendrycks2021mmlu} & \textbf{79.0\%} (5,648q) & \$0.00 & $-$9pp vs GPT-4o \\
GPQA Diamond (20q) & \textbf{90\%} & \$0.012 & $>$Claude 3.5 \\
RouterBench AIQ \cite{withmartian2024routerbench} & \textbf{0.6673} & \$0.00 & best published \\
m$\mu$Score & \textbf{10,607} & \$0.00 & \#1 all routers \\
Human oracle & \textbf{100\%} (200/200) & \$0.00 & 200 prod.\ queries \\
\bottomrule
\end{tabular}
\caption{Benchmark results. HumanEval achieved by local 9B model. GPQA 90\% driven by Tier~0 resolving physics constants, CS bounds, and biological facts deterministically.}
\label{tab:benchmarks}
\end{table}

\subsection{Local vs.\ Cloud-Cheap on Code (MultiPL-E, 180 problems)}

The qwen3-coder:30b local model \textbf{outperforms GPT-4o-mini class}
in 8 of 12 languages (74.4\% local vs.\ 64.4\% cloud). Local wins on
all compiled/statically-typed languages (C++, C\#, Go, Java: 80\%
local vs.\ 0--75\% cloud); cloud wins on dynamic scripting (PHP, R,
Ruby, Rust). \textbf{The local tier is not a compromise --- it is the
stronger model for the query categories it handles.}

\subsection{Routing Quality}

\begin{itemize}
  \item \textbf{RouterBench AIQ 0.6673} --- area under cost-quality
    convex hull, highest published value (prior best: KNN $\sim$0.65,
    RouteLLM BERT $\sim$0.62~\cite{ong2024routellm})
  \item \textbf{DeBERTa accuracy: 82.4\%} on held-out 517-example
    validation set (5,161 total training examples, 8~epochs CUDA,
    $\sim$22\,min)
  \item \textbf{Human oracle: 100\%} (200 production queries; 3
    initially disputed, subsequently conceded as correctly routed)
  \item \textbf{m$\mu$Score \#1:} quality $\times$
    $\sqrt{\text{cost\_efficiency}}$ $\times$ 10,000 = \textbf{10,607}
    vs.\ next-best 7,957 (bella)
\end{itemize}

\subsection{Tier Contribution (Ablation)}

\begin{table}[t]
\centering
\small
\begin{tabular}{lccc}
\toprule
\textbf{System} & \textbf{Free\%} & \textbf{Cost} & \textbf{Reduction} \\
\midrule
\textbf{Full system (Tiers 0--3)} & \textbf{87.5\%} & \textbf{\$27.87} & \textbf{87.3\%} \\
No Tier 0 & 78.5\% & \$29.46 & 86.6\% \\
No Tier 1 (cache) & 60.4\% & \$36.20 & 83.5\% \\
No Tier 2 (local) & 36.2\% & \$140.39 & 36.0\% \\
Cloud only & 0\% & \$219.26 & 0\% \\
\bottomrule
\end{tabular}
\caption{Tier ablation. Tier~2 removal multiplies cost 5$\times$. Tier~1 dominates latency. Tiers~0+1 resolve 38.9\% free on CPU-only hardware.}
\label{tab:ablation}
\end{table}

%% ============================================================
\section{Ecosystem Implications}

A distinctive property of this architecture: \textbf{costs fall as
usage grows}. Each resolved query contributes to cache density; each
cache hit reduces future marginal cost. This inverts the standard
inference cost curve.

Three behavioral properties emerge under zero-marginal-cost local
routing:

\textbf{Cost changes what users ask.} At \$0/query, usage shifts toward
exploratory queries users would suppress at \$0.01/query. The
15,363-query volume ($\sim$700 queries/day from a single developer) is
evidence of this behavioral shift. Benchmark distributions are biased
toward ``queries worth paying for,'' not queries people actually want to
ask.

\textbf{Privacy changes which questions get asked.} 87.5\% of queries
never leave the machine. Local-first routing creates a \textbf{privacy
floor by default} --- without configuration, consent dialogs, or explicit
action.

\textbf{Individual-scale deployment as a new regime.} Before capable
local inference, developers faced binary choice: cloud or nothing. The
\$600 equivalent cost (Sonnet-priced) for this window became \$27.87.
We propose \textit{individual-scale deployment} as an underexplored
research regime with properties (domain specialization, cache
compounding, privacy floor) invisible to population-scale benchmarks.

%% ============================================================
\section*{Limitations}
\label{sec:limitations}

\textbf{Single-user scope.} All 15,363 queries come from one developer
over 22 days. The DeBERTa classifier was trained and validated on the
same distribution. Generalization to multi-user, domain-specific
(legal, medical), or non-English deployments is not established.
Classifier accuracy varies substantially by category: Reasoning 87.5\%,
Coding 46.3\%.

\textbf{Hardware dependency.} Tier~2 requires $\geq$24\,GB VRAM GPU
(tested: RTX~5090 32\,GB). CPU-only deployments access Tiers~0+1
(38.9\% free), degrading gracefully to cloud for inference queries.

\textbf{Privacy residual.} The 12.5\% of queries reaching Tier~3 are
transmitted to Anthropic, OpenAI, or Google servers and subject to
their data retention policies. Users with strict data-locality
requirements should disable Tier~3 (\texttt{DISABLE\_CLOUD=1}).

\textbf{Benchmark sample sizes.} GPQA Diamond (20 questions) and GSM8K
(20 problems) samples are small; results should be interpreted as
directional rather than definitive.

%% ============================================================
\section{Conclusion}

muLLM demonstrates that a four-tier routing hierarchy --- deterministic
lookup, semantic cache, local inference, cloud fallback --- achieves
96.4\% cost reduction over a blended cloud baseline while maintaining
100\% HumanEval pass@1 and 79.0\% MMLU at \$0.00. The key insight:
production query complexity is not uniformly distributed. The majority
of queries are repetitive, pattern-matchable, or structurally simple,
and these can be resolved at zero marginal cost before any neural
inference is required. The system gets cheaper the more it is used.

%% ============================================================
\bibliography{mullm_references}
\bibliographystyle{acl_natbib}

%% ============================================================
\appendix

\section{m$\mu$PETS --- Pareto Efficiency Timing Score}

For routing systems, code quality and reasoning benchmarks measure the
underlying model rather than the router. m$\mu$PETS combines only the
three dimensions where routing decisions have direct causal impact:

\[
\text{m}\mu\text{PETS} = \frac{\text{routing\_accuracy} + \text{cost\_efficiency} + \text{ttft\_score}}{3}
\]

\begin{table}[h]
\centering
\small
\begin{tabular}{lcc}
\toprule
\textbf{Component} & \textbf{muLLM} & \textbf{Always-GPT-4} \\
\midrule
Routing accuracy & 0.824 & 1.00 \\
Cost efficiency & 0.987 & 0.00 \\
TTFT score & 0.85 & 0.30 \\
\midrule
\textbf{m$\mu$PETS} & \textbf{0.887} & \textbf{0.43} \\
\bottomrule
\end{tabular}
\caption{m$\mu$PETS scores. A system routing everything to GPT-4 scores 0.0 on cost efficiency. A cloud-only baseline (GPT-4o-mini) scores $\sim$0.55.}
\label{tab:mpets}
\end{table}

\textbf{AI assistance disclosure.} AI tools were used for software
development and manuscript drafting. AI tools are not listed as authors
per ACL policy.

\textbf{Ethics.} No human subjects were recruited. 87.5\% of queries
were processed entirely on-device. Cloud-routed queries (12.5\%) were
transmitted to providers under their standard API terms.

\end{document}
